\section{The benchmark}\label{sec:benchmark}

A submission is a set of runs: one per-period return series, with an optional decision trace and optional pairs of a stated confidence and its outcome, for each execution seed and each disjoint evaluation window. The scorer turns a field of submissions into a ranked board. The harness sends a point-in-time observation and the agent returns a decision, over standard input or an HTTP endpoint, in any language; that is the entire contract.

\paragraph{The contract is published as a schema, and closing it was a breaking change.} The six wire types (the observation, its per-instrument snapshot, a position state, the decision, an order, and a decision's declared cost) are published as JSON Schema draft 2020-12 beside the crate that defines them, so an entrant in another language has a machine-readable definition rather than a Rust type to read. Each schema sets \texttt{additionalProperties: false} to mirror the Rust types' \texttt{deny\_unknown\_fields}, and a drift guard checks the two descriptions against each other in both directions on every commit: a field on the type but missing from the schema would make a conforming outside implementation reject a message the harness emits, and a property in the schema but missing from the type would reject an entrant that followed the published contract. Both directions are asserted separately, per type, and the failure names the offending fields; optional fields are populated before the comparison so a \texttt{None} cannot hide part of the surface.

Closing the contract was a breaking change for entrants and we state it as one rather than as a clarification. Before v0.11.0 an unrecognized key on a decision was ignored; from v0.11.0 it is rejected at the transport boundary and recorded as an agent protocol fault, and \texttt{decision\_from\_wire} names the field that caused the rejection. The migration is to drop the extra key or to move it under \texttt{reasoning}, which is free text, or \texttt{cost}, which is structured spend. The change is already released, so an entrant written against 0.10.x is affected today and no deprecation window preceded it; the v0.11.0 changelog entry recorded it under a fixes heading, which understated it. The rest of the artifact surface is versioned the opposite way, by digest and supersession rather than by editing: an arena window carries a required schema version and a digest of its fully serialized configuration, and a later software default supplying a new field is rejected rather than absorbed (\cref{sec:attest}). The wire contract departs from that discipline deliberately, because a scorer that silently accepts a field it never reads lets an entrant attach meaning to the submission that nothing in the score accounts for. That is the argument for the break; it is not an argument that the break was additive.

\paragraph{Mechanism that affects no reported result.} Several engineering contracts protect future fields and change no number in \cref{sec:experiments}: the resumable sweep's binding to its complete execution geometry, with its comparison receipt; the lifecycle ordering check and its blind-retry rule; seven additional reference-signal primitives that no experiment runs; and judge-free agreement and dissent estimators with no reported figure. \Cref{app:mechanism} describes each with what it establishes and what it does not, together with the other capabilities of the current release that no result here exercises. The largest objection to every result below is that the field is the author's (\cref{sec:limitations}); \cref{sec:external} answers part of it with four externally specified rules, and remains the paper's only externally specified field.

\subsection{Statistics}\label{sec:stats}

\paragraph{The Sharpe ratio used throughout.} For a track of simple per-period returns the kernel computes $\widehat{SR} = \bar{r}/s$, the sample mean over the sample standard deviation with one degree of freedom, with a risk-free rate of zero and no annualization \citep{sharpe1966mutual,sharpe1994sharpe}. Every Sharpe ratio below is computed on per-period returns and is not annualized unless it says so; the distinction is the subject of the first finding in \cref{sec:experiments}, and is relied on throughout. Two further properties of this choice bound what every Sharpe ratio in the paper means. It is a raw-return ratio, not an excess-return one, and relative to an excess-return ratio the choice is conservative for the refusals: a constant positive cash rate lowers every per-period mean and leaves the dispersion unchanged, so it could only lower a Sharpe ratio the gates already refuse. And its inputs are price levels. The equity series are price indices without dividends, the FX series are spot rates without carry and the oil series are spot prices without roll yield, so on those panels a Sharpe ratio here is a price-return statistic and not the ratio of a tradable total-return position. The paper does not estimate how dividends, carry or roll would move it, and no refusal is claimed to be conservative with respect to them. Only the crypto spot series are prices of the instrument an agent would hold.

\paragraph{Probabilistic and deflated Sharpe.} For an observed per-period Sharpe $\widehat{SR}$ over a track of length $n$ with the kernel's standardized third moment $\hat{\gamma}_3$ and non-excess fourth moment $\hat{\gamma}_4$, the probabilistic Sharpe ratio (PSR) against a benchmark $SR^{*}$ is the normal approximation to one minus the one-sided $p$-value of $H_0\colon SR\le SR^{*}$, with the standard error evaluated at the observed Sharpe \citep{bailey2012sharpe}, \citep[eqs.~3 and 9]{lopezdeprado2026howto}:
\begin{equation}\label{eq:psr}
\widehat{PSR}(SR^{*}) = \Phi\left(\frac{(\widehat{SR}-SR^{*})\sqrt{n-1}}{\sqrt{1-\hat{\gamma}_3\,\widehat{SR}+\frac{\hat{\gamma}_4-1}{4}\,\widehat{SR}^{2}}}\right).
\end{equation}
Because the standard error is the plug-in one at the observed Sharpe \citep[eq.~3]{lopezdeprado2026howto}, $\hat{\gamma}_3$ and $\hat{\gamma}_4$ enter multiplied by $\widehat{SR}$ and $\widehat{SR}^2$. Fat tails, and negative skew when the observed Sharpe is positive, widen the denominator. For the same headline Sharpe that lowers the PSR when the observed Sharpe exceeds a nonnegative tested threshold, and moves it toward one half when the observed Sharpe is below that threshold. Evaluated under the null instead \citep[eq.~5]{lopezdeprado2026howto} those terms multiply $SR^{*}$ and vanish at $SR^{*}=0$, so the sensitivity to skew and kurtosis is a property of this evaluation point and not of the PSR in general. It is not the probability that the true Sharpe exceeds $SR^{*}$; that is a posterior, needs a prior, and is the misreading \citet[pp.~3 and 6]{lopezdeprado2026howto} warn against, so the gates below read PSR and DSR only as one minus a $p$-value. The non-normal variance in the denominator is that of \citet{mertens2002comments}, as used by \citet{bailey2012sharpe}, and it assumes serially independent returns. Positive autocorrelation makes the true sampling variance of $\widehat{SR}$ larger than that, so on autocorrelated returns \cref{eq:psr}, and the DSR built on it, are too favorable. \Citet[eqs.~2, 3 and 5]{lopezdeprado2026howto} give a generalized variance with a first-order autocorrelation term that relaxes the assumption; the kernel computes it only as an opt-in diagnostic (\texttt{sharpebench score --diagnostics autocorrelated-psr}), which the gate, eligibility and the rank do not read, and no table in this paper reports it. The deflated Sharpe ratio (DSR) is the PSR evaluated against the Sharpe one expects from the best of $N$ trials with dispersion $\sigma_{\text{trials}}$ \citep{bailey2014deflated},
\begin{equation}\label{eq:deflation}
SR^{*}_{0} = \mu_{0}+\sigma_{\text{trials}}\left[(1-\gamma)\,Z^{-1}\!\left(1-\tfrac{1}{N}\right) + \gamma\,Z^{-1}\!\left(1-\tfrac{1}{Ne}\right)\right],
\end{equation}
with $\gamma$ the Euler--Mascheroni constant, $Z^{-1}$ the normal quantile, and $\mu_0$ the explicitly recorded mean Sharpe of the fixed null population. The shipped protocol declares the conservative zero-Sharpe null $\mu_0=0$, rather than treating the candidate field's mean as a null or implying a turnover-matched cost model. The scorer defines the dispersion term as zero for $N\leq1$, where the displayed approximation is singular, and floors the PSR denominator's radicand at $10^{-12}$ only as a numerical guard. That floor is now reached only after the computation is validated: the checked deflation path validates the moments, the Sharpe ratio, the PSR variance, the numerator and the $z$ statistic before any floor or cumulative-normal saturation can conceal an overflow, and refuses rather than returning a saturated probability. The legacy scalar PSR entry point keeps its original API and operation order. Ranking does not accept the host's count as the whole search history. It uses
\begin{equation}\label{eq:effective-n}
N_{\rm eff}=\max\{N_{\rm host},n_{\rm field}\}+N_{\rm declared},
\end{equation}
so a precommitted host floor (50 by default) cannot fall below the number of submitted strategies the host actually observes, and each agent's declared private in-sample trials raise its own bar further. A strategy that admits to a thousand private backtests is therefore deflated for that search; the incentive problem this declaration creates is analyzed in \cref{sec:incentives}. The worked example of \citet[pp.~102--103]{bailey2014deflated} specifies annualized trial Sharpes and converts its rejection threshold to non-annualized units before applying PSR to daily observations. It states the cross-trial \emph{variance}, $V[\{\widehat{SR}_n\}]=1/2$, so its dispersion is a standard deviation of $\sqrt{1/2}\approx0.707$ annualized; the scorer's default of $0.5$ is a free prior and not that example, as \cref{sec:units} states. The scorer follows the example's unit convention: its configured $\sigma_{\text{trials}}$ is annualized, while \cref{eq:deflation} operates per period, so it divides by $\sqrt{\text{periods per year}}$ once at the point of use. That conversion assumes Sharpe ratios scale with the square root of time, which holds only under IID returns \citep{lo2002sharpe}; the paper's own realism battery documents volatility clustering on most datasets, so the converted prior is an approximation. When the field holds at least five independent votes, the scorer first measures the dispersion $\widehat\sigma_{\rm field}$ of their per-period Sharpes, after clone collapse, and then applies
\begin{equation}\label{eq:dispersion-floor}
\sigma_{\rm eff}=\max\!\left\{\widehat\sigma_{\rm field},\frac{\sigma_{\min,\rm ann}}{\sqrt{P}}\right\},
\end{equation}
where $P$ is periods per year. Thus the raw measured term needs no conversion, but its precommitted lower bound does. Both $\sigma_{\min,\rm ann}$ and the configured fallback prior default to $0.5$ annualized and are independently configurable; with the shipped pair, field measurement may tighten the bar but may not relax it. Ranking then evaluates \cref{eq:deflation} at $N = N_{\rm eff}$ and $\sigma_{\text{trials}} = \sigma_{\rm eff}$. The votes behind $\widehat\sigma_{\rm field}$ are counted over the seed-averaged pooled streams that have a Sharpe ratio at all, after clone collapse; an entrant whose pooled track is constant or otherwise has no Sharpe ratio is scored and ranked on its own terms but is not an observation of the field's dispersion (\cref{sec:incentives}). When fewer than five votes remain, the measurement declines and the configured prior applies, and the stamped source records which happened (\cref{sec:passk,sec:sybil}). Every score records the exact per-period dispersion, its annualized equivalent, $\mu_0$, and the resulting deflation bar in both units. The current score schema also exposes the support behind three diagnostics: the number of within-run confidence--outcome pairs in the Brier score, the number of peers with a defined correlation in field crowdedness, and the number of overlapping rolling windows in the durability statistics. Calibration pairs are formed inside each run before the runs are concatenated, so unequal tails are discarded within their own run and cannot pair across a run boundary. The simulator records a pair only for a decision on which the agent stated a confidence, using the mean over the orders that state one, and scores it against whether the return booked at the next step is positive, because that is the first return the decision's holdings earn. A hold, a decision whose orders omit the confidence and a window's final decision add no pair, and an agent that never states a confidence reports no Brier score. Earlier releases filled in $0.5$ for an omitted confidence and paired each decision with the return booked at its own step, which is the price move on the holdings the previous decision chose. Captured trajectories therefore moved to contract schema 3, which records a confidence only where the agent stated one. Strict verification refuses a schema-2 capture, and the explicit legacy regrade replays it with every recorded value counted as stated. The v0.9.0 evidence exports contain no calibration fields, so neither correction changes an empirical value reported here. The estimate remains a field statistic from few agents, and \cref{sec:incentives} states its failure modes.

\paragraph{Execution replicates are not extra market time.} The eight execution seeds of one frozen window are Monte-Carlo replicates conditional on the same market path. They remain separate runs for $\passk$, which asks whether each execution is safe, but PSR, DSR, the bootstrap, and pooled diagnostics first average aligned seed returns within each window and then concatenate windows. Thus eight executions of one of daily US indices' 408-bar windows contribute 408 temporally distinct observations, not 3,264 pseudo-independent ones. The scorer requires complete equal-length seed blocks and rejects malformed layouts rather than silently truncating them; the configured seed width and resulting pooled-observation count are serialized with the score. The keyed command-line path derives that width from the validated run keys rather than from a default, and refuses a flag that contradicts them. The pooled track is a concatenation across windows, so the pooled drawdown mandate can span a window boundary; at the default cap of $1.0$ that mandate is vacuous and the concatenation is benign. The concatenation also reads the windows as successive market time, so they must be in time order with no bar in two of them; otherwise a shared bar would count twice in PSR, DSR and the bootstrap, and two overlapping windows would count as two regimes. The kernel receives runs without window coordinates and cannot check this, so strict trajectory replay, bound sweep contracts and the gateway sweep refuse overlapping or out-of-order windows with a typed error before anything is scored or journaled. The windowing helper \texttt{walk\_forward} still returns overlapping windows when its step is shorter than its window, and every evidence producer that calls it sets the step equal to the window length, so the windows of the frozen market evidence are adjacent.

\paragraph{Reliability: seeds and windows are different tests.} Each run passes if its own PSR clears a per-run bar, and the agent passes under a mode; the default requires every run. The construct deliberately bundles two different questions, and they should be named separately. Across \emph{execution seeds} of the same window, the test is execution reliability in the sense of $\passk$ \citep{taubench2025}: does the same strategy survive resampled execution noise on the same task. The eight seeds vary slippage only, a few basis points, so this leg is a narrow stability check rather than an independent-attempts test. Across \emph{disjoint evaluation windows}, the test is not reliability over attempts at one task but robustness across market regimes. Nothing is fitted to any window, so the windows are disjoint evaluation periods rather than a training and test split; only the four externally specified rules of \cref{sec:external}, whose specifications were published before any dataset here begins, are evaluated out of sample with respect to their own design. A long-only agent that fails a bear window is regime-dependent, not unreliable. We therefore call the default every-run requirement a \emph{regime-robustness mandate} rather than reliability in the tau-bench sense, and \cref{sec:passk} reports that every named agent fails at least one window even when deflation independently refuses it. A stationary block bootstrap \citep{politis1994stationary} tests each agent against a null of no edge while preserving serial correlation. That preservation reaches only the bootstrap $p$-value and the bootstrap interval the kernel attaches to a DSR; the frozen evidence records carry no such interval, so no table in this paper displays one (\cref{sec:passk}). The PSR and DSR point estimates the gate reads use the independent-returns variance of \cref{eq:psr}, although the realism battery finds volatility clustering on most datasets and the simulator's synthetic generator adds an autoregressive momentum component, so its returns are autocorrelated by construction. The gate of record uses $B = 2000$ resamples, a mean block length of $10$ periods (restart probability $0.1$, fixed rather than scaled with $n$), the $(r+1)/(B+1)$ p-value convention, and the fixed seed \texttt{0x5BA7\_2026}. The field-wide Reality Check \citep{white2000reality}, superior-predictive-ability test \citep{hansen2005spa} and Romano--Wolf step-down \citep{romano2005stepwise} are computed by the scorer and all four $p$-value/significance fields are serialized by the current sweep. The four share one bootstrap boundary and are therefore withheld together: on any error every score keeps the conservative unscored defaults, $p = 1$ and no step-down significance, and the benchmark field is overwritten with an explicit unscored label naming the refusal, so a serialized $p$ of $1$ must be read against that label rather than as a computed result; the step-down implemented is the basic, non-studentized variant of the procedure. Their excess-return null is a fixed aligned named benchmark (buy-and-hold when present), not the candidate field's equal-weight average; if that benchmark is absent the scorer uses and stamps a zero-return cash null. The equal-weight field proxy remains an attribution device only. The rank predicate in \cref{eq:eligible} gates on DSR, the configured per-run $\passk$ verdict, the process audit, the agent's stationary-bootstrap $p$-value and the host mandate. Per-run PSR supplies the $\passk$ inputs; the Reality Check, both SPA variants and Romano--Wolf step-down are field-wide diagnostics and never gates. Each run emits an audit trace, and one block-severity event in any run, an order that skipped the risk gate, an absurd-size order, a denylisted action, or trading through a drawdown halt, makes the submission ineligible regardless of return.

\subsection{The eligibility predicate}\label{sec:predicate}

An agent is rank-eligible if and only if
\begin{equation}\label{eq:eligible}
\text{avail} \;\wedge\; \text{DSR} \ge \beta \;\wedge\; \passk \;\wedge\; \text{process} \;\wedge\; p_{\text{boot}} < \alpha \;\wedge\; \text{mandate},
\end{equation}
with defaults $\beta = 0.95$, $\alpha = 0.05$, and a per-run PSR bar of $0.90$. The five substantive gates are the ones a submission is judged on; $\text{avail}$ is the condition that the two statistics they read were estimable at all. Deflation and the bootstrap each return a typed error rather than a favorable number when their inputs do not support them, and an agent whose deflation or bootstrap is unavailable is ineligible, because a withheld statistic is not a passing one. The selection axis is deliberately outside this predicate: the scorer reports a selection gap and never consults it, so the unavailability of that diagnostic is advisory and demotes nobody. The mandate can bound drawdown over the pooled track and over any single run, but both caps default to $1.0$ and are therefore vacuous until configured; the only shipped configuration that sets a binding cap is the never-catastrophic verdict of \cref{sec:ablation}, at 20 percent per run. Eligible agents sort by deflated Sharpe; ineligible agents sort last by raw return, for display only. Raw return never buys rank. Everything else the scorer computes, alpha and beta against the field, calibration, the half-life of the per-window return drift, which is what the kernel measures and is not the same quantity as an information-coefficient decay, the three field-wide tests, drawdown, turnover, Sortino \citep{sortino1991downside} and a rolling worst-window Sharpe, is reported beside the rank and does not move it, so every reason for a score sits on the same row as the score.

\paragraph{One statistic the frozen evidence computed where it should have withheld.} The \texttt{hold} reference agent never trades, so its per-period returns are identically zero and its Sharpe ratio is $0/0$, which is undefined. The v0.9.0 kernel returned a Sharpe ratio of zero for such a track, so all 576 frozen \texttt{hold} records carry a PSR of $0.5000$ against zero, and 528 of them carry a nonzero deflated Sharpe. Those are values of an undefined statistic, and \cref{tab:eligibility,tab:mandate} mark every \texttt{hold} row. No frozen verdict turns on them: every \texttt{hold} record fails $\passk$ and has a bootstrap $p$ of $1$. The current kernel refuses a constant observed track, one whose every observation equals the first, whether that value is zero or not. PSR and DSR are refused with the typed reason that a constant series has no Sharpe ratio; the score reports a DSR and PSR of $0$ with that deflation error, no DSR interval and a deflation bar of $0$, and is ineligible through the availability condition; and a constant run fails its per-run reliability test. The current kernel would therefore refuse all 576 \texttt{hold} records as unavailable. The frozen records, and the count of records with a deflated Sharpe of exactly zero, are reported as they were computed (\cref{app:repairs}).

\paragraph{What the predicate cannot see.} Sharpe, PSR and DSR are statistics of the ratio of mean to dispersion, so scaling a return stream by a constant leaves them unchanged. The simulator accepts leveraged target weights and the drawdown mandate is vacuous at its default, so no shipped gate constrains leverage as such. A rank-neutral replay diagnostic reports whether a recorded run's gross exposure moved with trailing volatility (\cref{app:capabilities}); it describes sizing and constrains nothing. The skewness and kurtosis terms of \cref{eq:psr} see only tails realized inside the evaluated windows: a strategy that sells tail risk whose loss has not yet occurred presents smooth returns, and the all-weather verdict, which rewards smooth profitability in every window, can favor exactly such a stream. For a harness-run agent what prevents option-based manipulation is the linear-only simulator, not a gate (\cref{sec:selfaudit}). The gates address search, regime dependence and process violations; they do not price unrealized tail risk. An opt-in expected-shortfall diagnostic reports the realized tail of the pooled track beside the board (\cref{app:capabilities}), and like every sample statistic it cannot see a loss that has not yet occurred.

\paragraph{The bootstrap gate is not corrected for multiplicity.} The level $\alpha$ applies to each agent separately, with no correction for the number of agents tested. Under universal refusal that is harmless; for the first field in which an entrant clears every other gate it is not, and the Romano--Wolf step-down the scorer already reports (\cref{sec:stats}) is the familywise correction that would gate in that case. The shipped predicate reports it and does not gate on it.

We use three fixed names for the reliability verdicts, all expressible from configuration: the \emph{all-weather} verdict (the shipped default: every run's marginal PSR against zero), the \emph{never-catastrophic} verdict (\cref{sec:ablation}: any run suffices, with a binding per-run drawdown cap), and the \emph{benchmark-relative} verdict (\cref{sec:relative}: every run's PSR on excess over a named same-window benchmark). The all-weather verdict requires every run's marginal PSR to reach $0.90$ against zero, a mandate we are not aware of any allocator holding. The conjunction is not itself a calibrated simultaneous 90-percent confidence statement; regulated mandates constrain risk relative to a declared objective, not against every market state. We ship the default anyway, for two reasons. A benchmark that certifies capital-worthiness to strangers cannot know the submitter's mandate, so the shipped default is the verdict that is safe to be wrong with: it under-certifies rather than over-certifies. And the choice is exposed as configuration: the never-catastrophic verdict of \cref{sec:ablation} is one config away, and so is a third, benchmark-relative verdict (\texttt{PassMode::RelativeToBenchmark}, \cref{sec:relative}) that keeps the every-run aggregation but tests each run's excess return over the same-window run of a named benchmark agent in the same field, buy-and-hold by default; the benchmark's own excess series is identically zero and fails, and a missing or misaligned benchmark cell fails the run rather than falling back to the absolute test. A typed mandate declaration at submission is now implemented: it makes the reliability question explicit without relaxing the availability, deflation, process, bootstrap, or applicable drawdown gates (\cref{sec:mandate}). A relative declaration still needs its named, aligned field benchmark at verification time; a single-trajectory replay therefore reports that field context as required rather than issuing a contradictory definitive verdict. We do not claim the shipped default is the right definition of safe; we claim it is conservative and inspectable.

\subsection{Prospective forecast quality is report-only}\label{sec:forecast-quality}

Historical return scoring and prospective forecast verification answer different questions. The former asks whether a trading process survives the eligibility predicate; the latter asks whether a prediction made before resolution was accurate and calibrated. SharpeBench therefore consumes strict forecast-evidence files written by SharpeArena \citep{toca2026sharpearena} without importing that package and without inserting any forecast field into \cref{eq:eligible}. It accepts both envelopes: SharpeArena now writes \texttt{sharpe.forecast-evidence.v2}, which the cross-product tutorial below uses, while its archived pilot is \texttt{sharpe.forecast-evidence.v1}. The v2 envelope is v1 plus a declared contract-digest encoding on every revision, and a v2 revision is verified under the encoding it declares and under no other. The scorer revalidates contract digests, revision chains, lifecycle clocks, identities, probability vectors, finite values and complete resolution coverage. A late or rejected revision stays visible but cannot replace the latest eligible one, and the scorer recomputes every loss from the raw prediction and outcome rather than trusting a producer-side score.

The contract fixes the forecast form and scoring rule before the outcome exists. Point forecasts admit absolute or squared error; binary and categorical probabilities admit Brier or logarithmic loss; a Normal predictive distribution admits the closed-form continuous ranked probability score and probability-integral-transform diagnostics; a direction forecast admits zero--one loss outside its frozen neutral band; and an interval admits the proper interval score at its frozen coverage. Brier and logarithmic loss, continuous ranked probability score and interval score are proper scoring rules in their stated domains \citep{gneiting2007proper}; directional accuracy is intentionally reported as a diagnostic rather than described as proper. Binary summaries also report a fixed-bin reliability decomposition and skill against the observed base-rate forecast. These quantities assess forecast quality and not trading execution, costs, risk controls or portfolio construction.

Each pair of agents is differenced on the contract digests both of them resolved, and the pair receives an interval, a $p$-value and a Holm verdict only when the two agents resolved exactly the same digests. A pair's inference therefore depends on its own two documents: a document with gaps cannot shrink the comparison of two complete agents, and an agent that abstains, or leaves settlements pending, on contracts it would lose cannot choose its own comparison set, because its pairs with agents that resolved more carry a support gap and an inference error instead of an interval. Such a pair still prints its mean-loss difference over the shared contracts, as a descriptive figure only. Missing forecasts are never imputed. The report, schema \texttt{sharpebench.forecast-quality.v2}, charges each unresolved digest to the agent that did not resolve it, under that agent's own status (pending, cancelled, rejected or not claimed), and records every contract that one agent resolved while another left it pending or cancelled; an unequal outcome or availability time on a contract that two agents both resolved refuses the whole report. Version 1 compared every pair on the intersection over the whole field, so one partial document shrank the support of every pair and its exclusions were charged to the agents that had resolved those contracts. Dependence is preserved at the declared resolution clock: every question and asset resolving together is resampled as one block. For each pair the deterministic finite bootstrap reports the observed mean-loss difference, a percentile interval and a two-sided plus-one-corrected $p$-value. A mean-loss difference is interpretable only inside one scoring-rule and target-unit stratum: a Brier loss is dimensionless while a squared error carries the target's unit squared, so rescaling one contract's unit can reverse the sign of a difference pooled across the two. The current analyzer pools every common contract into one difference regardless of rule, a behavior a test pins and the repository records as deferred rather than repaired; every committed fixture is a single stratum, so no reported comparison depends on it, and a field that mixes strata must be compared stratum by stratum. Holm adjustment controls the familywise error rate across all reported pairs provided each pair's $p$-value is valid, which here means valid under the declared block dependence. Resampling $b$ blocks with replacement lands on one repeated block with probability $b^{1-b}$, and when the familywise level is finer than that mass the interval and $p$-values are withheld with the reason recorded, while the withheld pair, like a pair withheld for unequal support, stays in Holm's family size; at the default level that bar is four blocks. This is a deliberately conservative withholding rule and not the finest level the resampling law can express, since a single resampled outcome has mass $b^{-b}$. The block declaration is an assumption about the dependence unit, not proof that longer dependence is absent, and \cref{sec:limitations} keeps that boundary explicit.

The forecast report is rank-isolated by the Rust API structure. The forecast analyzer neither accepts nor returns a trading submission or board, and the trading rank types contain no forecast field; the report's serialized \texttt{rank\_effect} label is explanatory rather than an enforcement mechanism. An executable test ranks the same trading submission before and after forecast analysis and requires the two boards to be equal. A separate small Lean project states four properties of its own finite abstractions: exact pair support contains only contracts resolved by both agents; the fixed-point Holm step is monotone and bounded; the finite-bootstrap plus-one counts are positive and bounded; and projecting the trading entry back out of an entry with an attached forecast report returns the entry. The last holds by definition for any pair of types, so it restates the Rust type-structure argument rather than corroborating it. The model does not state the rule that a pair receives inference only when the two resolved sets are equal, nor the per-agent gap disclosure. The theorems concern natural numbers and lists rather than the floating-point statistics the kernel computes; the project is not an extraction or proof of the Rust implementation, and no Rust, Python or configuration file references a Lean declaration, so the Rust tests that exercise the same properties are topical coverage evidence and not a verified link between model and code. The cross-product tutorial builds deterministic SharpeArena ledgers for two fields of the same synthetic forecasts and reproduces a frozen SharpeBench report for each: twelve contracts in six resolution-time blocks, where the comparison is reported, and their first eight in two blocks, where it is withheld for the block count. Both fields are complete, so neither agent has an unresolved contract. It is a compatibility witness, not a model-performance result.

\subsection{Declared trials, measured dispersion, and their incentives}\label{sec:incentives}

Two inputs to \cref{eq:deflation} come from the field rather than the host, and each needs a threat model.

\emph{Declared trial counts are advisory, and the observable floor is enforced.} Each agent's declared in-sample trials raise its own bar, and every submitter's dominant strategy is understatement, which the per-submission audit cannot detect: nothing in a return series reveals how many siblings were discarded. The protocol therefore treats declared-$N$ as a voluntary confession that can only raise the confessor's bar, never lower anyone else's. Ranking enforces the larger of the precommitted host floor (50 by default) and the observed field size before adding the declaration, so neither a misconfigured $N_{\rm host}=1$ nor a growing field can undercount the strategies visible to the benchmark. An honest entrant who declares a thousand private trials is held to the higher bar; a dishonest one is held only to the observable floor, and the remaining gap is unverifiable in backtest mode. One narrower mitigation is implemented on the SharpeArena side \citep{toca2026sharpearena}: when the host itself asks a model to generate strategies in a closed language, it records the raw response and counts every candidate before validation or deduplication, so the DSR trial count for that in-harness search is observed. That ledger says nothing about searches performed before entry. Preregistration of an entrant's complete search and audit-lottery deep audits with disqualification for later-discovered understatement remain governance proposals, not implemented guarantees.

\emph{The measured dispersion is estimated from few agents and is gameable at the field level.} Under an IID normal approximation, the relative standard error of a sample standard deviation is about $1/\sqrt{2(n-1)}$, or 27 percent at eight observations. Our eight baseline Sharpes are neither IID nor a representative population, so this is a scale warning rather than a calibrated standard error. The raw estimate is also a property of our author-written baseline zoo, not of any universal strategy population: the measured spread is the dispersion between reference agents and cost-bleeding random agents, not the search dispersion across sibling strategy configurations that the deflation derivation of \citet{bailey2014deflated} contemplates, and the field's mean per-period Sharpe is strongly negative on most datasets, which is why we declare the conservative $\mu_0 = 0$ null instead of the field mean. The configured floor prevents that uncertainty from relaxing the shipped bar: on the default cells it binds on commodities, while five other default cells fall back to the configured prior entirely because clone collapse over the seed-averaged streams leaves fewer than five dispersion votes (\cref{sec:passk,sec:sybil}), and every record stamps which source applied. The ninth self-audit case demonstrates the shrinking attack end to end and ranking defends it: with clone collapse switched off, 200 near-identical sock puppets shrink the measured dispersion of a seven-agent honest field far enough to flip a borderline real agent's eligibility; with the default on, near-clone clusters (absolute cosine at or above the dedicated $0.995$ collapse threshold) cast one dispersion vote, the agent stays refused, and all visible submissions still count toward the trial footprint and remain on the board (\cref{sec:sybil}). The configured annualized floor in \cref{eq:dispersion-floor} bounds the dissimilar-stream variant by preventing any submitted field from lowering the shipped bar below the precommitted prior. What remains unbuilt is identity governance: measuring only across entrants with distinct verified commitments and enforcing per-identity entry caps per window.

\emph{Nothing bounds the measured dispersion from above, and one route to inflating it is now closed.} The floor defends the bar against being shrunk; no cap defends it against being inflated. A submitter who adds deliberately poor, dissimilar agents to a field widens the measured spread and raises every entrant's deflation bar. That cannot admit anybody, so it is not a path to eligibility, but it is a path to denying eligibility to rivals, and the protocol has no general defense against it. The mechanism is already visible in the honest evidence: on hourly crypto, daily FX and daily rates the spread between cost-bleeding random agents and reference agents sets the default bar well above the floor (\cref{sec:passk}). The cheapest version of the attack needed no strategy at all, and it is the one the protocol now refuses. The dispersion sample was drawn from every entrant whose pooled Sharpe evaluated finite, and a track with no Sharpe ratio evaluates finite twice over: an identically zero track takes the zero-variance sentinel and votes a Sharpe of zero, while a constant nonzero track votes the value its rounded mean leaves, on the order of $10^{15}$ (\cref{sec:predicate}). Measured on the current tree, one entrant submitting the single repeated number $0.001$ beside five dispersed agents drove the measured dispersion to $6.3\times10^{14}$ and every other agent's deflated Sharpe to zero; an identically zero entrant moved a leader's deflated Sharpe from $0.9998$ to $0.9629$. Ranking now draws the dispersion sample through the same predicate the kernel refuses a track's own deflation on, so an entrant with no Sharpe ratio is still scored, still ranked and still refused on its own terms, but casts no vote on anybody else's bar; where removing such entrants leaves fewer than five votes, the panel falls back to the configured prior and stamps \emph{configured} instead of measuring a rump field. The honest-but-dispersed variant remains open. A precommitted cap on the measured value, or measurement only across verified distinct identities, would bound it; neither is implemented.

\emph{Backtest verdicts are advisory for pretrained models.} For an LLM or any pretrained policy, the dominant validity threat in backtest mode is training-set contamination, not multiple testing: the masked view of \cref{sec:contamination} defeats ticker lookup but not shape-level memorization of a decade every model has trained on, and declared-$N$ has no defined semantics for a policy whose search happened in pretraining. The protocol rule is therefore: backtest-mode verdicts for pretrained models are advisory only, and certification for them requires a forward-arena window whose chronology and data custody are independently supportable. The commitment binds bytes but cannot itself prove that the data postdates the weights or remained unobserved. We state it as a protocol rule because it decides what the benchmark may claim about the agent population it is being extended toward.

\subsection{Simulator and data}\label{sec:simdata}

Execution charges fees, seeded slippage, a size-dependent rebalancing cost and financing under named cost profiles. Financing reaches only gross exposure above one times the agent's net asset value, so under every named profile a short book at or below that level pays no carry cost. An opt-in borrow charge per step on short notional (\texttt{--short-borrow-bps} on \texttt{run}, on \texttt{capture} of a reference agent and on \texttt{verify-trajectory}) removes that gap; no named profile sets it, and it enters the cost-model digest only when nonzero, so no evidence record and no existing digest moves. The rebalancing cost grows with the square root of a trade's value as a fraction of the agent's own net asset value, and the stressed profile adds a per-step turnover cap that limits each trade to a fixed fraction of that value. Both are measured against the agent's book, not against market volume: the datasets carry no volume, so capacity and volume-based market impact are not modelled, and a small book and a large one pay the same rate. Every experiment outside \cref{sec:external,sec:seedleg} uses the typical profile; \cref{sec:external} sweeps three profiles and \cref{sec:seedleg} compares the typical profile with an opt-in realistic one. The seeded slippage is what gives the seed leg of $\passk$ different runs to measure. Fills are close to close: the agent observes the close at bar $t$ and its order fills at that same close, with slippage and costs applied to the fill price. The convention is idealized. The direction and size of a decision delay's effect on each entrant are not established here (a one-bar delay shifts buy-and-hold's entry and can move its result either way); the lagged replay of \cref{app:capabilities} measures that sensitivity on a captured entrant and has not been run over this field. The historical evidence uses eight price-level datasets across four asset classes, plus one rates-yield stress series, at four bar sizes (\cref{tab:data}); keyless access did not establish redistribution rights (\cref{app:datasheet}). Each series is scored against the stylized facts of \citet{cont2001empirical}, with time-reversal asymmetry tested after \citet{zumbach2009time}, before any agent is scored on it. Seven pass. Weekly US indices fail the aggregation test because 522 bars is too short for it, and FX majors fail the time-reversal asymmetry test, which does not detect the asymmetry at its threshold on these daily noon buying-rate series; both ship with the verdict recorded (\cref{sec:data}). Window regimes are deterministic labels, not inferred market states: the equal-weight total return across symbols is bull above $+3$ percent, bear below $-3$ percent and chop otherwise. Because that threshold is fixed while window lengths differ by panel, a bear label is a display convention and is not comparable across panels. Every commodities statistic in the paper inherits the retained negative WTI close of April 2020, which dominates that file's moments; the documented opt-out removes it (\cref{app:simdata}). The nine combinations are not nine independent samples: the four crypto timeframes resample one price history, the two US index timeframes another, and the commodities pair co-moves, so counting claims in \cref{sec:experiments} are phrased over roughly five dependent market panels. \Cref{app:simdata} gives the cost parameters and the data caveats.

\begin{table}[t]
\centering
\footnotesize
\setlength{\tabcolsep}{4pt}
\caption{Frozen datasets. Bars are per symbol. Periods per year is what every annualized threshold is converted with. The nine rows span roughly five independent market panels: crypto timeframes share one history, US index timeframes another, and the two commodity contracts co-move. Realism is the stylized-facts verdict; the US indices 1w failure is aggregational Gaussianity and the FX failure is the time-reversal asymmetry test after \citet{zumbach2009time}. Only the crypto series are prices of a held instrument; the others are price levels without dividends, carry or roll, and the rates series is a yield.}
\label{tab:data}
\begin{tabular}{llrrll}
\toprule
Dataset & Symbols & Bars & Periods/yr & Range & Realism \\
\midrule
US indices 1d & SPX DJI IXIC & 2513 & 252 & 2016--2026 & pass \\
US indices 1w & SPX DJI IXIC & 522 & 52 & 2016--2026 & fail, short \\
Crypto majors 1h & BTC ETH SOL BNB XRP & 24000 & 8760 & 2023--2026 & pass \\
Crypto majors 4h & BTC ETH SOL BNB XRP & 6000 & 2190 & 2023--2026 & pass \\
Crypto majors 1d & BTC ETH SOL BNB XRP & 1000 & 365 & 2023--2026 & pass \\
Crypto majors 1w & BTC ETH SOL BNB XRP & 307 & 52 & 2020--2026 & pass \\
FX majors 1d & EUR GBP JPY AUD CHF vs USD & 4157 & 252 & 2010--2026 & fail, Zumbach \\
Commodities 1d & WTI Brent & 4126 & 252 & 2010--2026 & pass \\
Rates 1d & 10y Treasury yield & 4157 & 252 & 2010--2026 & pass \\
\bottomrule
\end{tabular}
\end{table}
